\begin{explanationbox}
\textbf{一句话理解}

内心不是随便取三点坐标的平均值，而是按\textbf{每个顶点所对的边长}来计算：对边越长，那个顶点前面的系数就越大。

\medskip
\textbf{先看具体数值：三边 $14,15,13$，三块面积怎样分？}
\VIncenterFigure{5}{.5}{14}{15}{13}{这一例 $r=4$：蓝块 $28$、橙块 $30$、紫块 $26$，恰好是三条底边 $14,15,13$ 的两倍。}
三条垂线同样长，都是 $r=4$，所以三个面积都等于对应底边乘以 $2$。
于是面积比 $28:30:26=14:15:13=a:b:c$。\textbf{共同的高被约去，留下的正是对边长。}

\textbf{为什么是对边，而不是旁边的边？}

$\vec{IA}$ 的系数 $a$ 来自小三角形 $IBC$ 的面积，它的底边是 $BC$；
$\vec{IB}$ 的系数 $b$ 来自 $ICA$ 的底边 $CA$；
$\vec{IC}$ 的系数 $c$ 来自 $IAB$ 的底边 $AB$。\textbf{每个顶点对应的是“不包含这个顶点的那条边”。}

\medskip
\textbf{另一个更简单的入口：先找角平分线上的点 $D$}

令 $D=AI\cap BC$，初中角平分线定理给出
\[
BD:DC=AB:AC=c:b.
\]
因此 $D$ 在 $BC$ 上的位置，可以由一个比例直接确定，而无须求角度；同理，$I$ 在 $AD$ 上的位置也可以用比例确定。\textbf{两次按比例分线段，就得到内心坐标。}
\VIncenterFigure{6}{.52}{14}{15}{13}{同一例中，先从 $B$ 向 $C$ 走全程的 $13/28$ 到 $D$，再从 $A$ 向 $D$ 走 $2/3$ 到 $I$。}

\medskip
\textbf{五种最实用的题目识别信号}
\begin{enumerate}
\item \textbf{已知三个顶点，求内心坐标：} 先算 $a,b,c$，再按对边乘坐标，不求角平分线方程。
\item \textbf{已知 $a\vec{PA}+b\vec{PB}+c\vec{PC}=\vec0$：} 看系数是否依次等于（或成比例于）对边长，直接判断 $P=I$。
\item \textbf{求内角平分线方向：} 用两边方向的单位向量相加；但它只给方向，不能代替求交点。
\item \textbf{一条角平分线交另一边：} 用 $BD:DC=c:b$，把一个未知交点变成已知比例。
\item \textbf{题目只涉及三个小三角形面积：} 优先考虑它们拥有相等的高，不必先引入向量。
\end{enumerate}

\textbf{什么时候不值得用？}

如果三角形是直角三角形、两直角边沿坐标轴，直接由内切圆半径找圆心可能更快；如果题目已经给出角平分线与边的交点比例，也可以直接用初中比例关系。\textbf{选择让未知量最少的办法。}

\medskip
\textbf{与重心的区别}

重心的三个顶点系数同为 $1$：$\vec{OG}=(\vec{OA}+\vec{OB}+\vec{OC})/3$；
内心的系数则是对边长 $a,b,c$。只有等边三角形中，这两个点才会重合；例5用同一张图比较它们的位置。
\end{explanationbox}
